\documentclass[11pt]{article}
\usepackage[T1]{fontenc}
\usepackage{lmodern,microtype,amsmath,amssymb,amsthm,booktabs,array}
\usepackage[margin=1in]{geometry}
\usepackage[authoryear,round]{natbib}
\usepackage[hidelinks]{hyperref}
\hypersetup{pdftitle={Sharp subsidies for truthful production and network cooperation},pdfauthor={Maxwell Flanagan}}
\newtheorem{theorem}{Theorem}[section]
\newtheorem{lemma}[theorem]{Lemma}
\newtheorem{proposition}[theorem]{Proposition}
\newtheorem{corollary}[theorem]{Corollary}
\theoremstyle{definition}\newtheorem{definition}[theorem]{Definition}
\theoremstyle{remark}\newtheorem{remark}[theorem]{Remark}
\newcommand{\R}{\mathbb R}
\newcommand{\E}{\mathbb E}
\newcommand{\Prob}{\mathbb P}
\newcommand{\core}{\mathcal C}
\newcommand{\pos}[1]{(#1)_+}
\newcommand{\ind}{\mathbf 1}
\title{Sharp subsidies for truthful production\\and network cooperation}
\author{Maxwell Flanagan\\[2pt]\small New York City, United States\\[2pt]\small Corresponding author: \texttt{maxwell.flanagan@gmail.com}}
\date{17 October 2026\\[2pt]\small Economic Theory candidate; not submitted}
\begin{document}
\maketitle
\begin{abstract}
Owners of complementary productive resources may gain by withholding capacity
when a sharing rule selects efficient core allocations. We ask how much
external subsidy suffices to retain every coalition's production guarantee
while making each fixed owner's payment nondecreasing in its reported
resource bundle. For any technology represented by the minimum of $K$ fixed
nonnegative linear forms, we construct a deterministic rule satisfying both
requirements at all nonnegative real report profiles. Its total payment is
at most $(1+o(1))T\log K/\log\log K$, where $T$ is grand-coalition production.
A matching lower bound holds for the worst case with at most $K$ structural
cuts even on simple directed acyclic flow networks with distinct strategic
edge owners. Thus the leading constant is one. The upper combines capped
supporting prices selected without the recipient's report with endpoint
marginal payments. The lower transports payments across unilateral report
changes and explicitly accounts for auxiliary owners. The incentive model
permits binding componentwise underreports and payment-only utility.
The bound is uniform in owner count and resource dimension, but $K$ may be
exponential in a compact description. Exact finite evaluation exists;
polynomial preprocessing is not asserted. The result isolates the cost of
linking coalition-covering allocations across reports in games whose
individual profiles already admit efficient core allocations. For every finite
monotone concave homogeneous technology we also obtain total-payment factors
$n\sqrt{2/(n+1)}$ and $1+2\sqrt d$ in fixed owner count $n$ and resource
dimension $d$. Their unrestricted optimal rates remain open.
\end{abstract}
\noindent\textbf{Keywords:} production game; mechanism design; capacity reporting;
coalitional stability; subsidy; flow game.\\
\textbf{JEL classification:} C71; D82; D47.

\section{Introduction}
Resource owners can jointly create value that none can produce alone. A
sharing rule must then address two distinct problems: retaining the coalition
and eliciting the resources available to it. At a given capacity profile,
competitive supporting prices can allocate the entire production value while
meeting every coalition's outside option. The difficulty is choosing these
allocations consistently across profiles. An owner who receives less after
revealing additional resources has an incentive to withhold them.

We measure the external budget needed to reconcile these requirements.
Owners, resource types and technology are fixed. Each owner reports a
nonnegative bundle and can commit no more than its true resources. Its utility
is its payment; resource use has no private cost. A single deterministic rule
must work at every nonnegative real profile, including zeros and ties. Payments
are nonnegative, cover every coalition's production value, and are
nondecreasing in the reporting owner's entire bundle with the other reports
held fixed. These assumptions describe capacity disclosure by fixed partners;
they do not model hidden opportunity costs or the choice of who joins.

The production function is a minimum of $K$ fixed nonnegative linear forms.
This includes finite nonnegative linear production programs and
single-commodity flow networks. A form can represent supporting resource
prices; in a network it is a structural cut. Grand production is denoted by
$T$. If total payment is at most $BT$, the external subsidy is at most
$(B-1)T$. Coverage includes the grand coalition, so total payments cannot
be below $T$ at positive-value profiles.

\subsection{Result and economic interpretation}
We prove that the optimal worst-case total-payment factor is
\[
 (1+o(1))\frac{\log K}{\log\log K}
\]
for technologies with at most $K$ forms. The same leading constant is necessary
even on simple directed acyclic networks with one strategic owner per edge.
Theorem~\ref{prod:upper} provides an explicit finite-$K$ bound and a mechanism.
Theorem~\ref{dag:sharp} supplies the matching asymptotic lower bound for the
worst case with at most $K$ cuts. Owner and resource counts are unrestricted
finite integers in both statements.

The lower bound means that no universal constant subsidy factor preserves
all coalition guarantees and truthful capacity disclosure throughout this
class. This is not caused by an empty core at a fixed profile. For any
minimizing form $w$, the efficient payments $w\cdot b_i$ cover every
coalition. Subsidy is required because payments must also satisfy incentive
inequalities between profiles. Our lower bound connects those inequalities
using actual unilateral changes by the original owners. All auxiliary edges
in its network construction belong to strategic owners and can have arbitrary
reports; no free infrastructure is introduced.

The upper bound selects supporting prices after randomly deleting whole
owners. A recipient's excluded price does not depend on its own report, and
its contribution is capped at grand production. Endpoint marginal payments
repair coalition coverage. Averaging over a short interval of deletion
probabilities controls the additional payment. Randomness defines a finite
expectation: the implemented rule is deterministic. The proof uses a union
bound over the fixed forms, which explains the technology-complexity parameter.

A small $K$ gives a budget guarantee independent of the number of owners.
However, $K$ can be exponential in an encoded network or production program.
Theorem~\ref{general:sqrt} gives separate bounds of order $\sqrt n$ and
$\sqrt d$, including technologies without a finite-form representation.
These bounds improve the project's earlier owner-count and dimension estimates,
but their unrestricted optimal rates remain open. No polynomial-time
preprocessing algorithm is asserted. It also makes no incentive claim for resource-use
costs, overreports, coordinated deviations, false names, or changes in the
owner population. These are separate economic questions.

\subsection{Related work and contribution boundary}
Dual-price core allocations are classical in linear production games
\citep{Owen,SametZemel}. The present contribution concerns selecting payments
across reports under a uniform budget bound. Fixed-game cost of stability
\citep{Bachrach} studies subsidy for coalitional stability; here efficient
cores already exist profile by profile.

The closest strategic comparison is Huang, Mei, Yang and Cao's
\emph{Mechanism Design in Max-Flows} (2025, inspected
\href{https://arxiv.org/html/2502.08248v1}{arXiv:2502.08248v1}, Sections 2 and 4).
It uses binding capacity underreports and payment utility, establishes an
obstruction to truthful efficient core selection, and studies efficient
truthful mechanisms without universal core membership, including its Shapley
rule (Proposition 1). We retain every coalition
inequality and permit subsidy. We do not claim their split/merge or
other-owner monotonicity guarantees.

\citet{Immorlica} impose cross-monotonicity when other owners join.
\citet{Branzei} define type-monotonic schemes on participating subgames of a
unit-resource market (Definition 1 and Theorems 5--6 of the final 2007 text).
Those comparisons change the population; ours changes one fixed owner's
reported bundle. No equivalence via strategic cloning is used.

\citet{ABV} do use own-resource monotonicity for a fixed alliance. Their
inspected author version defines a strictly concave network-utility model
and studies projection onto the full efficient core. Theorem 1 of that author version proves
own-resource monotonicity of the Shapley value, while noting that this value
can fail core membership. Its projection monotonicity evidence is computational,
not a universal proof. Online
Resource 1 supplies a strictly concave counterexample to a universal
extension of that evidence, with a version-specific comparison. This
auxiliary example is not an instance of our homogeneous main theorem.

Weighted multi-glove games and their total balancedness are established
\citep{MorettiNorde}. Their sharper special-case analysis, the connection to
the continuous Banzhaf functional, and additional mechanism comparisons appear
in Online Resource 1. Supporting prices, marginal contributions, concentration
bounds and weighted counting are established tools. The proposed contribution
is their combination into the joint all-profile coverage, incentive and sharp
budget guarantee, with a matching lower bound under fixed ownership. The
bounded source audit is not a certification of historical priority.

\subsection{Organization}
Section~\ref{prod:section} gives the model, mechanism and complete upper proof.
Section~\ref{general:section} gives complete owner-count and dimension proofs
for general concave homogeneous production. Section~\ref{dag:section} proves
the complexity lower bound and explains why naive subdivision would destroy
it. The conclusion separates the unrestricted parameter gaps from the
restricted mechanism barriers and computational questions. Online Resource 1 is an
indexed supplement containing the earlier hierarchy, general extensions and
exact obstructions; its repeated foundational proofs make it independently
readable. Online Resource 2 contains the scoped verification code and logs.

\section{General production: a sharp complexity bound}\label{prod:section}
There are $n$ fixed owners and $d$ fixed resource types. Owner $i$ reports
$b_i\in\R_+^d$; write $c=\sum_i b_i$. Fix a nonempty finite family
$V=\{w^1,\ldots,w^K\}\subseteq\R_+^d$, with a report-independent priority,
and define
\[
 f(c)=\min_{w\in V}w\cdot c,\qquad
 v_b(S)=f\left(\sum_{i\in S}b_i\right),\qquad T=f(c).
\]
A rule is eligible if its payments are nonnegative, cover every coalition,
and are componentwise nondecreasing in each owner's whole bundle, with all
other reports fixed. These conditions hold on all of $(\R_+^d)^n$.
They express truthfulness for binding underreports and payment-only utility;
there are no operating costs, overreport, false-name, population-change or
coordinated-report guarantees. A factor $C$ bounds \emph{total payment} by
$CT$; external subsidy is at most $(C-1)T$. The family $V$ remains fixed
across reports. Its size may be exponential in a compact description. All logarithms are natural.

\begin{theorem}[General sharp-order upper bound]\label{prod:upper}
For every integer $j\ge2$, there exists an eligible deterministic rule with factor
\begin{equation}\label{prod:constant}
 C_j(K)=j+\frac{\log K+2\log j+1}{\alpha_j},\qquad
 \alpha_j=\log\frac{j^3}{j^2+j-1}>0.
\end{equation}
With $L=\log K$ and
$j=\max\{2,\lceil L/[\log(2+L)]^2\rceil\}$, this factor is
$(1+o(1))\log K/\log\log K$, uniformly over owner counts, dimensions,
technologies and reports. Zero-value profiles receive zero; for $K=1$ the
constructed rule is efficient.
\end{theorem}

\paragraph{The mechanism.}
For $q\in(0,1)$ let independent owner bits satisfy $\Prob(B_i=0)=q$.
Choose the first minimizer of the \emph{uncapped} objective,
\[
 W(b,B)=\operatorname*{first\,argmin}_{w\in V}
               \sum_i B_i(w\cdot b_i),\qquad
 a_i(w,b)=\min\{w\cdot b_i,T\}.
\]
Only the payment contribution $a_i$ is capped. Put
\begin{align}
 z_i^q&=\E a_i(W,b),&
 p_i^q&=\E[a_i(W,b)\mid B_i=0],\label{prod:p}\\
 m_i^q&=\E_{B_{-i}}\left[f\left(b_i+\sum_{h\ne i}B_hb_h\right)
                         -f\left(\sum_{h\ne i}B_hb_h\right)\right],
 &r_i^q&=qp_i^q+(1-q)m_i^q.\label{prod:hybrid}
\end{align}
The final payment is the deterministic finite expectation
\begin{equation}\label{prod:payment}
 M_i(b)=j\int_0^{1/j}r_i^q(b)\,dq.
\end{equation}
These quantities are finite polynomials in $q$ at a fixed profile; no
continuity or generic-position assumption about the reports is required.

\begin{proof}[Proof of Theorem~\ref{prod:upper}]
Conditional on $B_i=0$, every score and the selected form are independent of
$b_i$. For that fixed nonnegative form, both $w\cdot b_i$ and $T$ are
nondecreasing in $b_i$, so their minimum is nondecreasing. In $m_i^q$ the
off term is independent of $b_i$ and the on term is nondecreasing. The
nonnegative fixed mixtures preserve this property and nonnegativity.

For any $w,S$, if some $i\in S$ has $w\cdot b_i\ge T$, then
$\sum_{i\in S}a_i(w,b)\ge T\ge v_b(S)$. Otherwise no term in $S$ is capped,
and the sum is $w\cdot\sum_{i\in S}b_i\ge v_b(S)$. Thus every capped form
covers coalitions. Fix $B_{-i}$ and let $w_1$ be the on minimizer. The off
optimum is at most the off objective of this same form, whence
\begin{equation}\label{prod:endpoint}
 f\left(b_i+\sum_{h\ne i}B_hb_h\right)-f\left(\sum_{h\ne i}B_hb_h\right)
 \ge w_1\cdot b_i=a_i(w_1,b).
\end{equation}
The equality holds because the contribution is a nonnegative summand of
the selected on value, which is at most $T$. Averaging proves $r_i^q\ge z_i^q$,
and hence coverage for $r^q$ and $M$.

For the budget define $g(q)=\E f(\sum_iB_ib_i)$, $Z(q)=\sum_i z_i^q$,
and $R(q)=\sum_i r_i^q$. Capping leaves every open contribution unchanged:
\[
 \sum_iB_i a_i(W,b)=f\left(\sum_iB_ib_i\right).
\]
Consequently $q\sum_i p_i^q=Z-g$. Differentiating the finite mask polynomial
gives $\sum_i m_i^q=-g'$, so $R=Z-g-(1-q)g'$. Since $g(0)=T$ and $g\ge0$,
for $0<\eta<1$ integration by parts gives the exact identity
\begin{equation}\label{prod:integral}
 \int_0^\eta R\,dq=\int_0^\eta Z\,dq+T-(1-\eta)g(\eta)
                         -2\int_0^\eta g\,dq
 \le\int_0^\eta Z\,dq+T.
\end{equation}
Both the sign and the factor two in this identity matter.

Suppose $T>0$. For each fixed form let
$d_i(w)=a_i(w,b)/T\in[0,1]$ and $\rho_w=\sum_i d_i(w)$.
The selected form satisfies $\sum_iB_id_i(W)\le1$. For $\lambda>0$ put
$\alpha=-\log(q+(1-q)e^{-\lambda})$. The convex function
$u\mapsto\log(q+(1-q)e^{-\lambda u})$ lies below its endpoint chord
$-\alpha u$ on $[0,1]$. Independence and Markov's inequality give
\[
 \Prob\left(\sum_iB_id_i(w)\le1\right)
 \le e^\lambda\prod_i(q+(1-q)e^{-\lambda d_i(w)})
 \le e^{\lambda-\alpha\rho_w}.
\]
A union over the fixed family and integration of the resulting tail yield
\begin{equation}\label{prod:chernoff}
 \Prob(\rho_W\ge x)\le\min\{1,K e^{\lambda-\alpha x}\},\qquad
 \frac{Z(q)}T\le\frac{\log K+\lambda+1}{\alpha}.
\end{equation}
Take $\eta=1/j$ and $\lambda=2\log j$. For $q\le1/j$,
$q+(1-q)e^{-\lambda}\le(j^2+j-1)/j^3$, so
\eqref{prod:integral}--\eqref{prod:chernoff} give \eqref{prod:constant}.
For the displayed choice of $j$, as $L\to\infty$,
$j\sim L/(\log L)^2$, $\alpha_j\sim\log L$, and
$j=o(L/\log L)$, proving the asserted asymptotic.

If $T=0$, all masked production values and capped contributions vanish.
If $K=1$, no contribution is capped, $p_i^q=m_i^q=w\cdot b_i$, and
$M_i=w\cdot b_i$. Finally each finite-mask term is
$q^{n-1-|S|}(1-q)^{|S|}[q a_i(W(b,1_S),b)+(1-q)(v_b(S\cup\{i\})-v_b(S))]$.
Its integral is a finite sum with rational coefficients when $j$ is integer.
This gives exact rational payments for rational input, without an efficiency
claim about enumeration of masks or forms.
\end{proof}

\paragraph{A worked payment calculation.}
Take two fixed owners, $f(c)=\min\{c_1,c_2\}$, the ordered forms
$V=((1,0),(0,1))$, and $j=2$. Owner 2 reports $b_2=(0,2)$ throughout;
compare owner 1's bundle $(3,0)$ with its binding underreport $(1,0)$.
The report space remains the full two-owner bundle domain. At either displayed
profile, write $T=f(b_1+b_2)$. When owner 1 is excluded, the first form is
selected for either value of $B_2$, including the all-zero tie, so $p_1^q=T$.
When owner 2 is excluded, its capped contribution is $T$ exactly when owner 1
is present, so $p_2^q=(1-q)T$. Both endpoint marginals equal $(1-q)T$.
For example, at the full report and mask $(B_1,B_2)=(0,1)$, owner 1's
raw contribution $3$ is capped at $T=2$. Thus
\[
 r_1^q=T(1-q+q^2),\qquad r_2^q=T(1-q),\qquad
 (M_1,M_2)=\left(\frac{5T}{6},\frac{3T}{4}\right).
\]
\begin{center}
\begin{tabular}{cccccc}
\toprule
Owner 1's report & $T$ & $M_1$ & $M_2$ & Total payment & Subsidy\\
\midrule
$(3,0)$ & $2$ & $5/3$ & $3/2$ & $19/6$ & $7/6$\\
$(1,0)$ & $1$ & $5/6$ & $3/4$ & $19/12$ & $7/12$\\
\bottomrule
\end{tabular}
\end{center}
The underreport reduces owner 1's payment. Singleton production values are
zero and total payment is $19T/12\ge T$, so every coalition is covered at
both profiles. These payments illustrate the specified rule, not a necessary
subsidy: the efficient two-owner Shapley alternative discussed in
Proposition~\ref{prod:uncapped} pays $(1,1)$ and $(1/2,1/2)$ respectively.

\begin{corollary}[Production LPs and networks]\label{prod:lp}
The theorem applies to
$f(c)=\max\{\pi\cdot x:x\ge0,Ax\le c\}$ for fixed $A\ge0$, with every
activity column nonzero and fixed positive profits $\pi$. Here $V$ can be
the vertices of $\{w\ge0:A^Tw\ge\pi\}$. For single-commodity directed flow,
one may instead use the characteristic vectors of the inclusion-minimal
structural source--sink edge cuts.
\end{corollary}
\begin{proof}
Each activity is bounded by a positive entry in its column, so every feasible
production set is compact at finite $c$. The dual is nonempty, pointed and
polyhedral. Strong duality and attainment at a vertex give the claimed finite
min-linear representation, including boundary capacities. Network max-flow/
min-cut gives the cut representation directly. Redundant nonminimal cuts may
be removed once and for all because capacities are nonnegative. All owners
and forms remain fixed across profiles. A positive-profit zero-resource
activity is excluded because it would make production unbounded.
\end{proof}
Dual core allocations are established ingredients~\citep{Owen,SametZemel}.
General dual vertices may be fractional; their count is not an ordinary
network-cut count. Constraints introducing capacity-independent dual offsets
are outside the homogeneous model here.

\begin{proposition}[Why the payment cap is essential]\label{prod:uncapped}
Removing the cap from the own-excluded contribution in
\eqref{prod:p} makes the uniform-$q$ hybrid unbounded even with two owners
and two supporting forms.
\end{proposition}
\begin{proof}
Take unit-profit activity columns $(1,1)$ and $(M,0)$, $M>1$, and two
fixed owners with bundles $b_1=(M,0)$ and $b_2=(0,1)$, so $c=(M,1)$. The two dual vertices are $(1,0)$ and $(1/M,1-1/M)$, and
$T=2-1/M$. If owner 1 is off and owner 2 on, $(1,0)$ uniquely minimizes the
raw score, giving uncapped contribution $M$. Thus $r_1^q\ge q(1-q)M$;
with $j=2$ its average is at least $M/6$, an unbounded ratio to $T$.
This obstructs only the uncapped rule. On the full two-owner bundle
domain, the Shapley payments
$P_i(b)=[f(b_i)+T-f(b_{-i})]/2$ are efficient and own-monotone.
They are nonnegative by monotonicity of $f$, and cover singleton coalitions
because $f(b_1+b_2)\ge f(b_1)+f(b_2)$. This last inequality follows by
minimizing nonnegative linear forms. Thus an efficient eligible alternative
exists without restricting each owner's report to a single resource.
\end{proof}

\section{Bounds in owner count and resource dimension}\label{general:section}
The number of supporting forms can be much larger than the economic description.
We now obtain bounds directly in the numbers of fixed owners and resources.
Here the production technology may be any finite-valued monotone concave
function $f:\R_+^d\to\R_+$ homogeneous of degree one; a finite min-linear
representation is not required. The report space and eligibility conditions
remain unchanged. In particular, there is no additional fairness floor.

\begin{theorem}[General owner and dimension bounds]\label{general:sqrt}
For every fixed $n,d\ge1$ and every such technology, there exists an eligible
deterministic rule on all nonnegative real report profiles with
\begin{equation}\label{general:combined}
 \sum_i M_i(b)\le
 \min\left\{d,\ n\sqrt{\frac{2}{n+1}},\ 1+2\sqrt d\right\}T.
\end{equation}
The individual bounds are uniform in the other parameter. They also hold when
$f$ has downward discontinuities at the boundary. More precisely, for every
fixed $0<a<1$ the supporting-price/endpoint rule averaged over $q\in[0,a]$
has both bounds
\begin{align}
 B_n(a)&=\frac{na}{2}+
 \frac{n[1-(1-a)^{n+1}]}{(n+1)a},\label{general:owner}\\
 D_d(a)&=\frac{1+2d[-a-\log(1-a)]}{a}.\label{general:dimension}
\end{align}
The owner expression remains valid at $a=1$ by polynomial integration.
No polynomial-time evaluation or preprocessing claim is made.
\end{theorem}

\subsection{Supporting prices on boundary faces}
For $x\ne0$, restrict $f$ to its coordinate face
$F_x=\{z\ge0:z_e=0\text{ whenever }x_e=0\}$.
At the relative-interior point $x$, finite concavity gives a supergradient $u$.
Monotonicity forces each active coordinate of $u$ to be nonnegative, and
homogeneity applied above and below $x$ gives $u\cdot x=f(x)$.
For a boundary point $z\in F_x$, apply the supporting inequality at
$(1-t)x+tz$, $0<t<1$, and the concavity lower bound on that segment to obtain
$u\cdot z\ge f(z)$. Thus boundary continuity is unnecessary.
Choose the unique minimum-norm vector in the nonempty closed convex set of
such nonnegative supports on $F_x$. Extend its linear form to
\[
 \ell_x(z)=\begin{cases}u\cdot z,&z\in F_x,\\+\infty,&z\notin F_x.
 \end{cases}
\]
Set $\ell_0(0)=0$ and $\ell_0(z)=+\infty$ for $z\ne0$.
This is one fixed selection depending only on $x$. For finite min-linear
technologies the original fixed-priority global minimizing form is an
alternative; the bounds below apply to either choice.

For any fixed mask $B\subseteq N$, put $x_B=\sum_{i\in B}b_i$ and
$A_i(B)=\min\{\ell_{x_B}(b_i),T\}$. This vector covers every coalition:
if any member's raw price reaches $T$, its capped payment alone suffices;
otherwise all relevant prices are finite on the support face and their sum
majorizes the coalition's worth. For $i\in B$,
\[
 f(x_B)-f(x_B-b_i)\ge \ell_{x_B}(b_i),
\]
because the support touches at $x_B$ and majorizes $f(x_B-b_i)$.
The marginal is at most $T$. Replace open owners' prices by these marginals
and keep $A_i(B)$ for closed owners. Coverage is preserved. A closed owner's
selected price excludes its own report, while its raw price and cap are
nondecreasing in that report. An open owner's marginal has a fixed off term
and a nondecreasing on term. This proves eligibility for every fixed mask
and every report-independent mixture of masks.

With independent missing probability $q$, retain the notation $p_i^q,m_i^q,r_i^q$
from \eqref{prod:p}--\eqref{prod:hybrid}, now using $\ell_{x_B}$.
At a fixed profile every capped mask payment is finite, so the expectations
are finite polynomials in $q$. Define $M_i^a=a^{-1}\int_0^a r_i^q\,dq$.
The endpoints have measure zero, including when $a=1$.
Writing $g(q)=\E f(x_B)$ and $C(q)=q\sum_i p_i^q$, differentiation of the
finite mask polynomial gives
\begin{equation}\label{general:endpoint}
 \sum_i r_i^q=C(q)-(1-q)g'(q),\qquad
 \int_0^a-(1-q)g'(q)\,dq
 =T-(1-a)g(a)-\int_0^a g(q)\,dq.
\end{equation}
No derivative of $f$ with respect to reports occurs here.

\subsection{Owner-count bound}
Every closed payment is at most $T$, whence $C(q)\le nqT$.
The all-open mask alone gives $g(q)\ge(1-q)^nT$.
Substituting both estimates in \eqref{general:endpoint}, integrating and
dividing by $a$ yields
\[
 \sum_i M_i^a\le\left(\frac{na}{2}+
 \frac{n[1-(1-a)^{n+1}]}{(n+1)a}\right)T.
\]
Dropping the bracketed power term gives
$B_n(a)\le na/2+n/[(n+1)a]$. The fixed choice
$a=\sqrt{2/(n+1)}\le1$ equalizes these terms and proves the owner bound
in \eqref{general:combined}. For $n=1$ the exact expression is $B_1(a)=1$.
The former $O(n/\log n)$ upper bound recorded in the supplement is therefore
superseded by $(\sqrt2+o(1))\sqrt n$.

\subsection{Resource-dimension bound}
The following scalar identity controls contributions from closed owners
without counting owners or supporting forms.
\begin{lemma}[Bernoulli reciprocal bound]\label{general:reciprocal}
Let $u_i\ge0$, independently retain each $i$ with probability $p=1-q$, and
write $S_B=\sum_{i\in B}u_i$. Define $\min\{u_i/S_B,1\}$ to be zero when
$u_i=0$ and one when $u_i>0=S_B$. Then
\[
 \E\sum_{i\notin B}\min\{u_i/S_B,1\}\le\frac{2q}{1-q}.
\]
\end{lemma}
\begin{proof}
With zero-over-zero interpreted as zero,
$\min\{u/S,1\}\le2u/(u+S)$. Conditioning on the own bit gives
\[
 \E\sum_{i\notin B}\frac{u_i}{u_i+S_B}
 =\frac qp\E\sum_{i\in B}\frac{u_i}{S_B}
 =\frac qp\Prob(S_B>0)\le\frac qp.
\]
In the first equality both conditioned ratios are
$u_i/(u_i+\sum_{j\ne i}B_ju_j)$; the other bits have the same law.
\end{proof}

Suppose $T>0$. For every mask and owner,
\begin{equation}\label{general:coordinate}
 \frac{A_i(B)}T\le
 \sum_{e=1}^d\min\left\{\frac{b_{ie}}{(x_B)_e},1\right\},
\end{equation}
with the lemma's conventions. If $b_i$ has a positive coordinate outside
$F_{x_B}$, the right side is at least one. Otherwise the selected nonnegative
support satisfies $\sum_eu_e(x_B)_e=f(x_B)\le T$, so each active term obeys
$u_eb_{ie}/T\le b_{ie}/(x_B)_e$. Apply
$\min\{\sum_e z_e,1\}\le\sum_e\min\{z_e,1\}$.
The case $x_B=b_i=0$ is immediate. The same argument applies to a finite
global minimizing form, since off-face prices remain capped by $T$.

Sum \eqref{general:coordinate} over closed owners and use
Lemma~\ref{general:reciprocal} separately in each resource coordinate:
$C(q)\le 2dqT/(1-q)$. The endpoint integral in
\eqref{general:endpoint} is at most $T$. Since
$\int_0^a q/(1-q)\,dq=-a-\log(1-a)$, this proves
\eqref{general:dimension}. Also
\[
 -a-\log(1-a)=\sum_{k\ge2}\frac{a^k}{k}
 \le\frac{a^2}{2(1-a)},\qquad
 D_d(a)\le\frac1a+\frac{da}{1-a}.
\]
Choosing $a=1/(1+\sqrt d)$ gives $1+2\sqrt d$.
At $T=0$, monotonicity makes all coalition values, caps and endpoint
marginals vanish, completing these cases.

\subsection{Baseline, selection and remaining gap}
The grand marginal $U_i=f(c)-f(c-b_i)$ is eligible: it is own-monotone and
dominates the support allocation at $c$. Put
$\alpha_i=\max_{e:c_e>0}b_{ie}/c_e$ when $T>0$.
Then $c-b_i\ge(1-\alpha_i)c$, so $U_i\le\alpha_iT$ and
$\sum_i\alpha_i\le d$. This proves the $dT$ baseline.
Choosing an entire rule using the fixed public $n,d$ before reports proves
\eqref{general:combined}; no pointwise minimum of payment vectors is used.
For a finite $K$ representation, the complexity bound in
Theorem~\ref{prod:upper} is another available choice.

The scalar leading constant two is already sharp as $q\downarrow0$:
two equal positive weights give expected closed ratio exactly $2q$.
This does not prove optimality of the dimension constant for the full mechanism.
The unrestricted owner and dimension rates remain open between the earlier
logarithmic-over-logarithmic lower bounds and the square-root uppers above.
The lower bounds and their parameter translations are proved in Online
Resource 1, theorems titled ``Owner-count range'' and ``Concave dimension gap''.
Online Resource 1 gives matching square-root growth obstructions only for
specified fixed-mask, common-support and endpoint architectures, including
a maximum-coalition-marginal cap variant. Those restrictions are essential;
they are not lower bounds for arbitrary eligible payment rules.

Square-root payment bounds also arise in static cost-of-stability theory:
Bachrach et al. (2009, arXiv:0907.4385v1, Theorem 3.5) and the published
\citet[Theorem 7]{Bachrach} give a total-funding factor $\sqrt n$ for
superadditive games without the cross-report incentive requirement. Every profile in our
model already admits an efficient support allocation, so its static cost of
stability is zero. The new bounds concern one rule linking those profiles.

\section{A matching lower bound on simple networks}\label{dag:section}
A network has one distinct fixed strategic owner per edge, and coalition
worth is its maximum source--sink flow using only its owned edges. Structural
cuts are inclusion-minimal edge sets disconnecting source from sink; their
number depends on the fixed graph, not its reported capacities. Define
$\mathcal B_{\rm dag}(K)$ as the supremum, over simple finite DAGs having at
most $K$ structural cuts, of the best all-profile eligible total-payment
factor on that graph. The analogous quantity for finite nonnegative
min-linear technologies with at most $K$ forms is $\mathcal B_{\rm prod}(K)$;
owner counts and dimensions are unrestricted finite integers.

\begin{theorem}[Sharp first-order worst-case factors]\label{dag:sharp}
As $K\to\infty$,
\[
 \mathcal B_{\rm dag}(K)=(1+o(1))\frac{\log K}{\log\log K},\qquad
 \mathcal B_{\rm prod}(K)=(1+o(1))\frac{\log K}{\log\log K}.
\]
The lower holds along an explicit sequence of simple DAGs at their exact cut
counts, and for the worst case with \emph{at most} $K$ cuts. It does not
assert an exact-count padding theorem for every integer $K$.
\end{theorem}

\begin{proof}
The upper follows from Theorem~\ref{prod:upper}. We give the simple-DAG lower,
which also supplies the production lower. Fix integers $d_0\ge1$, $m>d_0$,
$R\ge1$, and $k\ge2$. (Here $d_0$ is hierarchy height, not resource dimension.)
One module with terminals $a,b$ has vertices $L_1,\ldots,L_R$,
$u_1,\ldots,u_m$, $v_1,\ldots,v_m$, $W_1,\ldots,W_R$ and edges
\[
 aL_s,\quad L_su_j,\quad u_jv_j,\quad v_jW_t,\quad W_tb
 \quad(1\le s,t\le R,\ 1\le j\le m).
\]
Take $k$ copies in series by identifying consecutive terminals. Each middle
edge $u_jv_j$ has report in $\{1,m\}$; all auxiliary edges have report
$H=m^2$. Every edge, including every auxiliary edge, is a strategic original
owner whose full report domain is nonnegative. Only the finite profiles
used in this proof keep auxiliary reports fixed.

A type-$h$ module has $h$ high middle edges, $0\le h\le d_0$. For any subset
$J$ of middle edges, choose one left fan and one right fan, and include the
source/final fan edges plus the incident edges supporting exactly $J$.
This actual coalition has module flow $\sum_{j\in J}c_j$: the middle cut
bounds it above, and the selected branches carry it because their total is
at most $m^2$. In particular the full graph has
\begin{equation}\label{dag:value}
 T=m+(m-1)\min_{1\le l\le k}H_l.
\end{equation}
Choosing the two fans independently uniformly, any particular auxiliary
edge is included with probability at most $1/R$.

Symmetrize an arbitrary eligible factor-$B$ rule over module permutations
and branch permutations, relabeling each branch's incident auxiliary edges.
These are permutations of the fixed owner set preserving every coalition's
flow at all profiles: module flows combine by their minimum, so even module
permutations preserve coalition values. Averaging preserves all required
inequalities. At a histogram $\nu=(\nu_0,\ldots,\nu_{d_0})$ of module types,
write $x_h(\nu),z_h(\nu)$ for low and high variable-owner payments, and
\[
 Y(\nu)=\sum_h\nu_h[(m-h)x_h(\nu)+hz_h(\nu)],\qquad
 Q(\nu)=\text{total auxiliary-owner payment}.
\]
No auxiliary-payment symmetry is assumed; $Y+Q\le BT$.

First choose one low variable in each module, supported by random fans.
Averaged coalition coverage and nonnegativity give
$\sum_h\nu_hx_h+Q/R\ge1$. Next put $q_0=m-d_0+1$ and choose $q_0$ low
variables in each type below $d_0$, and one high variable in each type
$d_0$. These exist because $m>d_0$. The resulting coalition has worth at
least $q_0$, giving
$q_0\sum_{h<d_0}\nu_hx_h+\nu_{d_0}z_{d_0}+Q/R\ge q_0$.
Multiply the first inequality by $m-d_0$ and the second by $d_0$. With
$A=(d_0+1)m-d_0^2$,
\begin{equation}\label{dag:coalitions}
 A\le A\sum_{h<d_0}\nu_hx_h+
 \nu_{d_0}[(m-d_0)x_{d_0}+d_0z_{d_0}]+(m/R)Q.
\end{equation}

Fix $\alpha>0$, set $\lambda=A/\alpha$,
$Z=\sum_{h=0}^{d_0}\lambda^h/h!$, and $p_h=\lambda^h/(h!Z)$.
Average profiles with independent module types of probabilities $p_h$;
let $\mu(\nu)$ be their multinomial histogram weights. Raising one actual
low variable owner from $1$ to $m$ gives
$x_h(\nu)\le z_{h+1}(\nu')$ at
$\nu'=\nu-e_h+e_{h+1}$. Every other report, including every auxiliary report,
is unchanged. The exact identity
\[
 \mu(\nu)\nu_h=(p_h/p_{h+1})\mu(\nu')\nu'_{h+1},\qquad
 A p_h/p_{h+1}=\alpha(h+1)
\]
transports only this changing owner's payment. Averaging
\eqref{dag:coalitions} therefore yields
\begin{align}
 A&\le\alpha\E\sum_{h=1}^{d_0}h\nu_hz_h+
 \E\nu_{d_0}[(m-d_0)x_{d_0}+d_0z_{d_0}]+(m/R)\E Q\notag\\
 &\le(1+\alpha)\E Y+(m/R)\E Q
 \le\max\{1+\alpha,m/R\}B\E T.\label{dag:transport}
\end{align}
The last step uses the disjoint budgets $Y,Q$; their coefficients are not
added. By the tail-sum formula and $p_0=1/Z$,
\begin{equation}\label{dag:mean}
 \E T=m+(m-1)\sum_{r=1}^{d_0}\left(\sum_{h=r}^{d_0}p_h\right)^k
 \le m[1+d_0 e^{-k/Z}].
\end{equation}

Let $C_{m,R}$ be the number of structural cuts of one module. Any minimal
cut is the outgoing cut of its reachable source side after deletion, so
$1\le C_{m,R}\le2^{2R+2m}$. A minimal cut in a series composition lies
entirely in one module: otherwise each module retains a path, or an edge
removed elsewhere is redundant. Hence the full cut count is exactly
$K=kC_{m,R}$.
For $d_0\ge2$ choose
\[
 g=\lceil\sqrt{\log(d_0+1)}\rceil,\quad
 m=d_0g,\quad R=m,\quad \alpha=1/g,\quad
 \ell=\lceil\log(d_0+1)\rceil,\quad k=\lceil2\ell Z\rceil.
\]
The separate multiplier $\ell$, rather than $g$, ensures the tail error
vanishes. Equations~\eqref{dag:transport}--\eqref{dag:mean} give
\[
 B\ge\frac{d_0+1-d_0/g}{(1+1/g)(1+d_0e^{-2\ell})}=(1-o(1))d_0.
\]
Here $\lambda=(d_0+1)d_0g^2-d_0^2g>d_0$, so
$\lambda^{d_0}/d_0!\le Z\le(d_0+1)\lambda^{d_0}/d_0!$.
Stirling's estimate implies
$\log k=(1+o(1))d_0\log d_0$; since $m=o(d_0\log d_0)$,
\[
 \log k\le\log K\le\log k+4m\log2=(1+o(1))d_0\log d_0.
\]
This proves the lower along the sequence $K_{d_0}$. For any fixed
$0<\varepsilon<1$ and large external cut budget $K$, take
$d_0=\lfloor(1-\varepsilon)\log K/\log\log K\rfloor$.
Then $K_{d_0}\le K$ eventually, yielding lower coefficient $1-\varepsilon$.
Letting $\varepsilon$ decrease to zero proves the at-most-$K$ assertion.
\end{proof}

\begin{proposition}[A subdivision that destroys the lower bound]\label{dag:subdivision}
Fully subdividing every edge of any finite directed network into two edges
owned separately produces a network with an efficient eligible rule.
\end{proposition}
\begin{proof}
For pair reports $(a_e,b_e)$ set $d_e=\min(a_e,b_e)$. Choose the original
network's fixed-priority minimum-cut core allocation $Q(d)$, which pays
$d_e$ on the selected cut and zero elsewhere. Pay $Q_e(d)$ to the higher
reporting member of pair $e$, split equally at equality, and pay the lower
member zero. Holding the mate's report fixed, the payment is zero below it,
half the fixed $Q_e$ at equality and the fixed $Q_e$ above it: in the last
two cases all virtual capacities are fixed. Thus it is own-monotone, even
if $Q$ is not. For a coalition $S$, let $S^*$ contain the original edges
whose two members both lie in $S$. Its value is exactly $v_d(S^*)$, and
$p(S)\ge Q(S^*)\ge v_d(S^*)$. Total payment is $T$. This shows why the naive subdivision reduction fails and motivates the
explicit auxiliary-budget accounting in the shared-fan construction.
\end{proof}


\section{Implications, limitations and further questions}
The result identifies technology complexity as one sharp measure of the
budget required to combine truthful disclosure with coalition coverage.
For every fixed finite representation, the upper holds simultaneously over
all profiles and all owner bundles. The lower uses an increasing family of
fixed technologies; it does not change the graph when an owner changes its
report. Its asymptotic statement concerns a worst case with at most $K$ cuts,
not an exact-count construction for each integer $K$.

Two questions remain distinct from the sharp $K$ theorem. First,
Theorem~\ref{general:sqrt} bounds the factor by
$\min\{d,n\sqrt{2/(n+1)},1+2\sqrt d\}$, while the unrestricted lower bounds
remain logarithmic over logarithmic in $n$ and $d$. The supplement's matching
square-root barriers apply to specified mask/support architectures; they do
not close these unrestricted gaps. A better general rule must escape those
restrictions. Second, rational finite technologies permit exact evaluation
by finite enumeration, and the supplement develops finite sampling with
all-profile certification. The certification preprocessing can be exponential.
A polynomial preprocessing guarantee or a matching computational obstruction
would be a separate contribution.

The main theorem deliberately isolates capacity disclosure. Private operating
costs, endogenous participation, and coordinated reports require additional
incentive constraints. A resource-sharing application should justify the
payment-only assumption and its feasible deviations before interpreting the
bound as a policy prescription.

\paragraph{Supplementary material.}
Online Resource 1 (\texttt{ESM\_1.pdf}): indexed proofs of supplementary
hierarchies, extensions and mechanism obstructions, including the original
research lineage. Online Resource 2 (\texttt{ESM\_2.zip}): scoped exact
verification programs, execution logs, source hashes and reproduction
instructions. Finite checks corroborate bounded identities and instances;
the universal and asymptotic statements depend on the proofs above.

\section*{Statements and Declarations}
\paragraph{Funding and competing interests.}
The author received no funding or other support for this work and has no
competing interests to declare.
\paragraph{Data and code availability.}
No empirical dataset is used in the main theorem. Online Resource 2 contains exact finite checks and their inputs, with
version hashes and execution records. No public repository deposit has been
made; the code and inputs are supplied with this manuscript.
\paragraph{AI assistance and responsibility.}
OpenAI Codex assisted the original research, mathematical exploration,
drafting, code and analysis. Anthropic Claude assisted subsequent internal
review. The author personally reviewed the manuscripts, proofs, code and
results. The historical Codex model version was not recorded. Internal model
review is not independent human peer review or formal proof verification.
No human specialist endorsement or comprehensive historical-priority
clearance is claimed.

\begin{thebibliography}{99}
\bibitem[Amigo et al.(2016)]{ABV}
Amigo I, Belzarena P, Vaton S (2016) Revenue sharing in network utility
maximization problems. NETNOMICS 17:255--284.
\url{https://doi.org/10.1007/s11066-016-9112-4}.
Comparison uses the author manuscript at \url{https://hal.science/hal-01400498v1}.
\bibitem[Bachrach et al.(2018)]{Bachrach}
Bachrach Y, Elkind E, Malizia E, Meir R, Pasechnik D, Rosenschein JS,
Rothe J, Zuckerman M (2018) Bounds on the cost of stabilizing a cooperative
game. Journal of Artificial Intelligence Research 63:987--1023.
\url{https://doi.org/10.1613/jair.1.11270}.
\bibitem[Br\^anzei et al.(2007)]{Branzei}
Br\^anzei R, Solymosi T, Tijs S (2007) Type monotonic allocation schemes
for a class of market games. TOP 15:78--88.
\url{https://doi.org/10.1007/s11750-007-0002-7}.
\bibitem[Immorlica et al.(2008)]{Immorlica}
Immorlica N, Mahdian M, Mirrokni VS (2008) Limitations of cross-monotonic
cost-sharing schemes. ACM Transactions on Algorithms 4(2), article 24.
\url{https://doi.org/10.1145/1361192.1361201}.
\bibitem[Moretti and Norde(2021)]{MorettiNorde}
Moretti S, Norde H (2021) A note on weighted multi-glove games.
Social Choice and Welfare 57:721--732. \url{https://doi.org/10.1007/s00355-021-01337-8}.
\bibitem[Owen(1975)]{Owen}
Owen G (1975) On the core of linear production games.
Mathematical Programming 9:358--370.
\url{https://doi.org/10.1007/BF01681356}.
\bibitem[Samet and Zemel(1984)]{SametZemel}
Samet D, Zemel E (1984) On the core and dual set of linear programming games.
Mathematics of Operations Research 9(2):309--316.
\url{https://doi.org/10.1287/moor.9.2.309}.
\end{thebibliography}
\end{document}
